\section{A fixed function from an independent set}
\label{ana:section}

The soundness argument begins by turning a large independent set into
one fixed Lipschitz function. We need a variance lower bound on fresh
random inputs and control of how much the function changes when one
position changes. Set
\begin{equation}\label{src:variance-parameters}
 b_0=\frac1{2m},\qquad
 \nu=\frac{b_0a^2}{4}=\frac1{2^7m^3}.
\end{equation}
Thus the independent-set density we must rule out is \(2b_0\).
First discard the low-norm vertices. All vertices
whose vectors have norm at most \(t\) form a clique: for any two of
them the triangle inequality gives
\(\norm{S_v+S_w}\leq2t\). An independent set therefore contains at
most one such vertex. Since
\(n\geq p>1/b_0\), an independent set of size at least
\(n/m=2b_0n\) contains more than \(b_0n\) vertices of norm greater
than \(t\). The following lemma turns these remaining vectors into
the required function.

\begin{lemma}[Separation, variance, and oscillation]
\label{ana:function}
Suppose an independent set contains at least \(b_0n\) vertices with
\(\norm{S_v}>t\). There is a fixed nonempty finite set
\(A_*\subset\R^{\mathcal P}\) such that
\[
 \Pr_{v\in V(G)}[S_v\in A_*]\geq b_0,
 \qquad
 \norm{x+y}>2t\quad(x,y\in A_*),
\]
where the vertex is uniform. For a nonempty finite set \(B\), write
\(\dist_*(z,B)=\min_{b\in B}\norm{z-b}\). The function
\begin{equation}
 F(z)=\frac12\bigl(\dist_*(z,-A_*)-\dist_*(z,A_*)\bigr)
 \label{ana:distance-function}
\end{equation}
is odd and \(1\)-Lipschitz for \(\norm{\cdot}\), and \(F(z)>t\)
on \(A_*\).

Let \(\mathcal E=\binom{[d]}2\). Draw independent uniform occurrence
seeds \(\boldsymbol c=(c_e)_{e\in\mathcal E}\in[M]^{\mathcal E}\)
and independent weight blocks \(s_j\) uniform on \([-1,1]^r\),
independently of the seeds. Put
\begin{equation}
 f(\boldsymbol c,\boldsymbol s)
    =F\left(\sum_{j=1}^d z(i_j(\boldsymbol c),s_j)\right).
 \label{ana:f}
\end{equation}
This bounded function satisfies
\begin{equation}
 \E_{\boldsymbol s}f(\boldsymbol c,\boldsymbol s)=0
       \quad\hbox{for every }\boldsymbol c,
 \qquad
 \E f^2\geq\nu d.
 \label{ana:variance}
\end{equation}
On the weight cubes it has the following three oscillation bounds:
\begin{enumerate}[label=\textup{(\roman*)}]
 \item Replacing one weight block \(s_j\), with all other data fixed,
       changes \(f\) by at most \(2\).
 \item Replacing one seed \(c_e\), with all other data fixed,
       changes \(f\) by at most \(4\).
 \item Fix \(k\in[d]\). Simultaneously replacing any seeds incident
       to \(k\), while keeping all nonincident seeds, all weights,
       and the question received at every opposite position
       \(j\ne k\) from \(\{j,k\}\) unchanged, changes \(f\)
       by at most \(2\).
\end{enumerate}
The set \(A_*\), the function \(F\), and its formula are fixed from
an independent set of the entire graph before the fresh seeds above
are drawn. A seed enters \(f\) only through its endpoint questions.
\end{lemma}

\begin{proof}
Take \(A_*\) to be the set of sum vectors of the stipulated high-norm
vertices in the independent set. The number of those vertices gives
the probability bound even if several vertices represent the same
vector. For distinct \(x,y\in A_*\), their representing vertices
are distinct and nonadjacent, so \Cref{graph:edge} gives
\(\norm{x+y}>2t\). If \(x=y\), the same inequality follows from
\(\norm{2x}=2\norm{x}>2t\).

Distance to a nonempty finite set is the minimum of the distances to
its points. The triangle inequality makes each distance function
\(1\)-Lipschitz; the half-difference in
\eqref{ana:distance-function} is therefore \(1\)-Lipschitz.
Symmetry of the norm gives
\(\dist_*(-z,-A_*)=\dist_*(z,A_*)\), proving oddness. For
\(z\in A_*\), its distance to \(A_*\) is zero, while
\[
 \dist_*(z,-A_*)=\min_{x\in A_*}\norm{z+x}>2t.
\]
Finiteness preserves the strict inequality, so \(F(z)>t\).

For fixed seeds, negating all weights preserves their joint law and
negates \(f\). This proves the first assertion in
\eqref{ana:variance}. Also \(F(0)=0\), and the increment bound
\(\norm{z(i,s)}\leq\|s\|_\infty\) implies \(|f|\leq d\)
on the sampling space.

Round every continuous entry to the midpoint of its grid interval.
The resulting entries are independent uniform elements of
\(\mathcal T\); interval-boundary choices have probability zero.
The change in one increment has norm at most \(1/p\), by linearity
in its weights and the increment bound. Thus rounding changes the
argument of \(F\), and hence its value, by at most \(d/p<t/2\).
The unchanged seeds together with the rounded weights have exactly
the uniform vertex law.
With probability at least \(b_0\), their sum belongs to \(A_*\),
and on that event the unrounded value of \(f\) exceeds \(t/2\).
Consequently
\[
 \E f^2\geq b_0t^2/4=b_0a^2d/4=\nu d.
\]

Every increment on a weight cube has norm at most one. Replacing
one weight block changes one increment by norm at most two, proving
\textup{(i)}. A single seed affects only its two endpoint queries;
replacing it changes at most two increments, each by norm at most
two, proving \textup{(ii)}. In \textup{(iii)}, each position other
than \(k\) keeps every question in its query, so only the increment
at \(k\) may change. The same bound of two proves the claim.
All these arguments use the single fixed function \(F\) in
\eqref{ana:distance-function}; it is not chosen again for a sampled
seed tuple.
\end{proof}
